\documentclass[11pt,a4paper]{article}
\pdfoutput=1 % if your are submitting a pdflatex (i.e. if you have
             % images in pdf, png or jpg format)
\usepackage[colorlinks=true, linkcolor=black!50!blue, urlcolor=blue, citecolor=blue, anchorcolor=blue]{hyperref}
\usepackage[font=small,labelfont=bf,margin=0mm,labelsep=period,tableposition=top]{caption}
\usepackage[a4paper,top=3cm,bottom=2.5cm,left=2.5cm,right=2.5cm,bindingoffset=0mm]{geometry}
\usepackage[T1]{fontenc} % if needed
%\usepackage{url,cite}

\usepackage[utf8]{inputenc}
\usepackage[english]{babel}
\usepackage{pifont}
\pdfoutput=1
\usepackage{graphicx,mathtools}
\usepackage{epsfig,cite}
\usepackage{amssymb}
\usepackage{amsmath}
%\usepackage{dsfont}
%\usepackage{multirow}
\usepackage{color}
\usepackage{amstext,alltt,setspace}
\usepackage{amsbsy}
%\usepackage{comment}
%\usepackage{fullpage}
\usepackage{array}
\usepackage{booktabs, makecell}
\usepackage{hyperref}
\usepackage{slashed}
\usepackage{url}
\usepackage{geometry}
\geometry{a4paper, textwidth=6.5in, textheight=9.5in, 
 marginparsep=7pt, marginparwidth=.6in,headheight=13.6pt}

\usepackage{amssymb,amsthm,cancel,hyperref,graphicx,xcolor}
\usepackage{picinpar,graphicx,xy}
\usepackage[subnum]{cases}
\usepackage{booktabs}
\usepackage{mathrsfs}
\usepackage{amsfonts}
\usepackage{latexsym}
\usepackage{booktabs,graphicx,hyperref,epsfig,multirow}
\usepackage{bm}
\usepackage{comment}
\setlength{\unitlength}{1mm}
\usepackage{tikz}
\usetikzlibrary{arrows.meta, positioning, fit, calc}
\usepackage{adjustbox}
\input{incl_shortcuts}

% commenting colors
\newcommand{\Eva}[1]{{\color{magenta}{EG: #1}}}
\newcommand{\ramon}[1]{{\color{red}{[RW: #1]}}}

% Only show up to subsection within TOC
\setcounter{tocdepth}{2}
\begin{document}
\begin{flushleft}
%%%%%%%%%%%%%%%%%%%%%%%%%%%%%%%%%%%%%%%%%%%%%%%%
\begin{figure}[h]
\includegraphics[width=.2\textwidth]{tiflab_logo.png}
\end{figure}
%%%%%%%%%%%%%%%%%%%%%%%%%%%%%%%%%%%%%%%%%%%%%%%%%
\end{flushleft}
\vspace{-5.0cm}
\begin{flushright}
TIF-UNIMI-2026-5 
\end{flushright}
\vspace{2.0cm}

\begin{center}{\Large \textbf{{
Interpreting Parton Distributions with Shapley Values
}}}\end{center}
\vspace{1.3cm}
\begin{center}
Rapha{\"e}l Bonnet-Guerrini$^{1,2}$,
Stefano Carrazza$^{2,3}$ 
Dakshansh Chawda$^{2,3}$,\\ 
Stefano Forte$^{2,3}$, 
Eva Groenendijk$^{2,3}$, 
Vincenzo Piuri$^{1}$, Ramon Winterhalder$^{2,3}$
\end{center}

\begin{center}
{}$^1$Dipartimento di Informatica, Universit\`a di Milano, 
Via Celoria 18, I-20133 Milano, Italy\\
{}$^2$TIF Lab, Dipartimento di Fisica, Universit\`a di Milano,  
and\\
{}$^3$INFN, Sezione di Milano, Via Celoria 16, I-20133 Milano, Italy

\vspace{0.9cm}


%%%%%%%%%%%%%%%%%%%%%%%%%%%%%%%%%%%%%%%%%%%%%%%%
%%%%%%%%%%%%%%%%%%% ABSTRACT %%%%%%%%%%%%%%%%%%%
%%%%%%%%%%%%%%%%%%%%%%%%%%%%%%%%%%%%%%%%%%%%%%%%

{\bf \large Abstract}
\end{center}

% We show that Shapley values can be used to trace the relationship between the behavior of individual parton distributions (PDFs) and  the  theory predictions obtained using them in the computation of high-energy observables.
We show that Shapley values can be used to trace how individual parton distributions (PDFs) shape the theory predictions for high-energy observables computed from them.
When determining PDFs from data, this provides a tool for assessing the impact of data on PDFs, while in the computation of collider observables, it provides a tool for assessing the impact of PDF uncertainties.
The Shapley value is computed by treating the regression of PDFs from data as a cooperative game. 
The PDFs are the players, and the reward is the likelihood ($\chi^2$) that characterizes the agreement between data and the predictions obtained from a given PDF, with theory and methodology held fixed. 
The method recovers known expectations about which data constrain which PDFs in a global fit, while placing them on a more quantitative footing. 
In some cases, it reveals unexpected results that admit a clear interpretation upon closer inspection. 
Though generally applicable, the method is particularly useful when PDFs are determined with a black-box AI model rather than a physically inspired parametrization. 
We demonstrate this by assessing the homogeneity of the data subsets used for the $K$-fold cross-validation in the NNPDF4.0 hyperoptimization.



\clearpage

\tableofcontents 



\clearpage
%%%%%%%%%%%%%%%%%%%%%%%%%%%%%%%%
\section{Introduction}
\label{sec:intro}

Parton distribution functions (PDFs) enter the computation of every process at a hadron collider, and how precisely they are known directly shapes the reliability of theory predictions.
This is a central challenge for precision particle phenomenology, with two distinct facets.
On the one hand, PDF uncertainties can hamper the search for new physics by masking the signal~\cite{Cole:2026eex}. 
On the other hand, since PDFs are themselves determined from data, understanding how the data constrain them is essential to optimize their extraction.  
This is nontrivial because PDFs are only indirectly related to the data, through perturbative QCD evolution and then convolution with a partonic hard cross section.


Several tools have been developed to study the sensitivity of PDFs to the data used to determine them~\cite{Wang:2018heo,Jing:2023isu,PDF4LHCWorkingGroup:2022cjn}. All of them assume that PDFs are found by maximum-likelihood determination of the parameters of a fixed functional form, and study the dependence of the likelihood on the parameters in the vicinity of its stationary point.
This setup does not directly apply when PDFs are determined with neural networks in place of a fixed functional form -- an approach proposed long ago~\cite{Forte:2002fg} and now central to the NNPDF methodology, which produces state-of-the-art sets of parton distributions~\cite{NNPDF:2021njg}.
In this case, the impact of the data on the PDFs is also obscured by the intrinsic opacity of the neural network, whose high-dimensional, non-linear structure makes it difficult to interpret directly~\cite{lipton2016_blackbox}.
Explainable AI (XAI) methods have been suggested in this context as a way to classify PDFs by their sensitivity to individual datapoints~\cite{Kriesten:2024are}. 
However, many standard XAI techniques that relate features to the behavior of the underlying model~\cite{ribeiro2016whyitrustyou, shrikumar2019learningimportantfeaturespropagating,binder2016layerwiserelevancepropagationneural} are not straightforwardly applicable here: the input is a one-dimensional momentum fraction $x$ rather than a high-dimensional feature vector with clear semantic interpretation, and the quality of the output can only be assessed after propagating the PDFs through the full theory pipeline.

In this work, we address this problem by adapting a specific XAI technique: Shapley values~\cite{shapley_value1953}, a cooperative game-theoretic concept that assigns each player their average marginal contribution across all possible coalitions.
Shapley values are model-agnostic and require no assumption on the internal structure of the system under study, which makes them naturally suited to our black-box pipeline. 
A common implementation, the SHAP framework~\cite{lundberg2017unifiedapproachinterpretingmodel}, has been shown to be potentially problematic in high-energy physics applications when correlations are important~\cite{Vent:2025ddm}, and this is specifically  the case for PDFs.
The reason is that SHAP relies on the so-called additivity assumption, namely that features are approximately independent. We tackle this issue by setting up the problem in such a way that the space of features is low-dimensional, which allows us to compute the Shapley values exactly while accounting for all correlations.

Our implementation of the Shapley value views individual PDF flavors as the players of the game. 
The reward of the game is the likelihood ($\chi^2$), computed as follows: starting from a set of PDFs determined with the NNPDF4.0 neural-network model~\cite{NNPDF:2021njg}, we introduce a controlled Gaussian perturbation in a localized range of the momentum fraction $x$, applied to one PDF at a time, for various choices of PDF basis.
Evaluating the Shapley value on different subsets of the global dataset used to determine the PDFs then allows us to disentangle the impact of each subset on each PDF and $x$ range. 
We show that the results largely reproduce standard lore on the impact of data on PDFs (as discussed in references such as~\cite{Forte:2010dt,Gao:2017yyd}), though in some cases they also reveal less obvious aspects of the connection between PDF features and data. 
As an application, we use the method to analyze the design of the folds in the default $K$-fold hyperoptimization of the NNPDF methodology. 

The paper is structured as follows. In Sec.~\ref{sec:PDF-nutshell}, we review the PDF-to-observable pipeline in the NNPDF machine-learning approach, and discuss its black-box nature. 
In Sec.~\ref{sec:shapley-values}, we introduce Shapley values, discuss the limitations of existing frameworks, and describe our exact implementation in PDF space, including the perturbation design and uncertainty quantification. 
In Sec.~\ref{sec:results}, we demonstrate the method on NNPDF4.0, scanning over $x$ regions and individual dataset families. 
We conclude with a summary and outlook in Sec.~\ref{sec:outlook}.


%%%%%%%%%%%%%%%%%%%%%%%%%%%%%%%%
\section{From PDFs to observables and vice-versa}
\label{sec:PDF-nutshell}

To make our treatment self-contained, we provide a brief overview of the NNPDF approach for determining PDFs, focusing specifically on the NNPDF4.0 implementation and methodology. The threefold goal is to establish a clear notation, make the paper accessible to readers with a computer science background, and highlight how the neural network model behaves as a black box, which we aim to elucidate using our XAI tool. We refer the reader to the NNPDF4.0 paper~\cite{NNPDF:2021njg} and its associated code and documentation~\cite{NNPDF:2021uiq} for a more detailed treatment. We first discuss how predictions for physical observables are computed from input PDFs, and then how PDFs are determined from data by solving the corresponding inverse problem in the NNPDF methodology.

%----------------------------------------
\subsection{Factorization and evolution}
\label{sec:factel}

PDFs are determined by performing a regression of theory predictions to experimental measurements of physical observables.
The factorization theorem expresses these predictions as convolutions of PDFs with perturbatively calculable partonic cross sections.
The PDFs characterize the quark and gluon content of the incoming hadrons, while the partonic cross sections describe the scattering of the corresponding initial-state quarks and gluons, collectively denoted as partons.
For processes with two incoming hadrons, this takes the form
%
\begin{equation}
    \mathcal{O}(Q^2) = \sum_{i,j} \int \dd x_{1}\, \dd x_{2}\;
    f_{i}(x_{1}, \muF^2)\,
    f_{j}(x_{2}, \muF^2)\,
    \hat{\sigma}_{ij}\!\left(x_{1}, x_{2},  \frac{Q^2}{\muF^2}\right),
    \label{eq:factorization_theorem}
\end{equation}
%
where $\mathcal{O}(Q^2)$ is a physical observable, here the inclusive cross section, depending on a hard scale $Q^2$, $f_i(x,\muF)$ is the PDF for flavor $i$ at momentum fraction $x$ and factorization scale $\muF$, $\hat{\sigma}_{ij}$ is the perturbatively calculable partonic cross section, and $\muF$ is a factorization scale, of which the observable does not depend, as the dependence cancels between the PDF and the hard cross section, up to perturbatively subleading terms. 
To improve perturbative convergence, the scale $\muF$ is chosen of the order of the physical scale $Q$, so that no large scale ratios arise. After imposing momentum conservation, the integral over $x_i$ takes the form of a convolution of the momentum fraction with some dimensionless ratio of physical scales on which both the observable and the partonic cross section depend. 
For lepton--hadron scattering, only one PDF is present, and the double convolution reduces to a single one. 

The scale dependence of the PDF is fixed by evolution equations of the form
%
\begin{equation}
\label{eq:evolution}
\frac{\partial f_i(x,\mu^2)}{\partial \ln \mu^2}
=
\sum_{j=1}^{2N_f+1} \int_x^1\frac{\dd z}{z}
P_{ij}\!\left(\alphaS(\mu^2),\frac{x}{z}\right)f_j(z,\mu^2),
\end{equation}
%
where the sum runs over the gluon and the $N_f$ quark and antiquark flavors with mass $m^2_i\lesssim \mu^2$, and suitable matching conditions when passing through the production threshold for heavy quarks.
The evolution kernels (splitting functions) in Eq.~(\ref{eq:evolution}) admit a perturbative expansion in powers of the strong coupling
\begin{equation}\label{eq:pertp}
P_{ij}\!\left(\alphaS(\mu^2),x\right)= \frac{\alphaS(\mu^2)}{2\pi} P^{(0)}_{ij}\!\left(x\right)+ 
\left(\frac{\alphaS(\mu^2)}{2\pi}\right)^2 P^{(1)}_{ij}\!\left(x\right)+\dots,
\end{equation} 
and give the
probability for of parton $j$ to emit a parton $i$ that carries a fraction $x$ of its
momentum.

Hence, even if PDFs must be evaluated at a variety of different scales $\mu$ to obtain accurate theory predictions, they can all be expressed in terms of a set of basis PDFs evaluated at a reference scale $\mu=Q_0$. 
By choosing a suitable (scale-independent) basis of linear combinations of PDFs, the evolution equations decouple into a pair of coupled equations, and $2N_f$ decoupled equations. We refer to this basis as evolution basis. The two coupled equations couple the gluon and the singlet combination, defined as
%
\begin{equation}\label{eq:singlet}
\Sigma(x,\mu^2)=\sum_{i=1}^{N_f} q^+_i(x,\mu^2), 
\end{equation}
%
where
%
\begin{equation}\label{eq:qpm}
q_i^\pm= q\pm \bar q,
\end{equation}
%
and $q_i$ and $\bar q_i$ are the quark and antiquark PDFs, respectively. The other decoupled equations evolve specific linear combination of quark and antiquark distributions, and have at most three independent evolution kernels (only one at leading and two at next-to-leading perturbative order), so several equivalent choices of evolution basis are possible.

At all scales, PDFs satisfy a set of integral constraints, namely the momentum sum rule
%
\begin{equation}
\label{eq:mom_sumrule}
    \int_0^1 \dd x\, x\left(g(x,\mu^2)+ \Sigma(x,\mu^2)\right)=1,
\end{equation}
%
and the valence sum rules
%
\begin{equation}
\label{eq:val_sumrule}
    \int_0^1 \dd x\, u^-(x,\mu^2) = 2\int_0^1 \dd x\, d^-(x,\mu^2)=2,\qquad  \int_0^1 \dd x\, s^-(x,\mu^2) = \int_0^1 \dd x\, c^-(x,\mu^2)=0,
\end{equation}
%
where $u$, $d$, $s$ and $c$ denotes the up, down, strange and charm quark PDFs, respectively.

%----------------------------------------
\subsection{PDF parametrization and determination}
\label{sec:pardet}

In the NNPDF approach the PDFs, viewed as a function of $x$ but at a fixed scale $Q_0$, are parametrized by a preprocessed neural network
%
\begin{equation}
\label{eq:PDF}
  x f_i(x, Q^2_0) = A_i\, x^{1-a_i} 
  \bigl(\text{NN}_i(x)-\text{NN}_i(1)\bigr),
\end{equation}
%
where $a_i$ are preprocessing exponents, $A_i$ are normalization constants that may be fixed to enforce the sum rules Eqs.~(\ref{eq:mom_sumrule}-\ref{eq:val_sumrule}), and the subtraction ensures that the PDFs vanish at the kinematic endpoint $x=1$. g

The neural networks are trained by minimizing a loss function, which takes the form of a negative log-likelihood describing the agreement between theory and data:
%
\begin{equation}
  L_{\text{NNPDF}} = \chi^2(f,\mathcal{D})= \sum_{n,m}^{N_{\text{dat}}}
  \left[ \mathcal{O}_{n}(f) - d_{n} \right] \Sigma_{nm}^{-1}
  \left[ \mathcal{O}_{m}(f) - d_{m} \right],
  \label{eq:loss_simp_pdf}
\end{equation}
%
where $\mathcal{O}_{n}(f)$ denotes the theoretical prediction for data point $n$ computed using a given PDF set $f$, $d_{n}$ is the corresponding experimental measurement of the full dataset $\mathcal{D} = \{d_1,\dots,d_{N_{\text{dat}}}\}$, and $\Sigma_{nm}$ is a covariance matrix that encodes both statistical and systematic experimental uncertainties as well as theory uncertainties originating from approximations in the theory predictions.  

The loss minimization is stopped through cross-validation, by dividing the data into training and validation samples and optimizing the agreement with the validation data. PDF uncertainties are determined by first, generating a Gaussian distribution of pseudodata replicas $\mathcal{D}^k$ whose mean and  covariance over the replica sample reproduce the original data and correlated uncertainties, and then, determining a best-fit PDF $f^k$ for each pseudodata replica, thereby leading to a distribution of best-fit neural networks whose mean and covariance represent the final determination of PDFs and associate correlated uncertainties. The training and validation partition is determined randomly for each replica, and the preprocessing exponents are varied randomly in a self-consistently determined interval that ensures that their choice does not bias the final result. Both the architecture and other features of the neural network (such as the initialization or the activation function) and the minimization algorithm and its features (such as the maximum gradient or the stopping patience) are determined through a $K$-fold hyperparameter optimization procedure, performed by dividing the data in four folds, excluding one fold at a time from the training process, and optimizing the agreement with the excluded fold.
Note that whereas in the NNPDF approach PDF uncertainties are natively represented  through Monte Carlo replicas, commonly available tools~\cite{Watt:2012tq} enable the conversion into the Monte Carlo replica representation of PDF uncertainties represented as multi-Gaussian in parameter space. The Monte Carlo representation is accordingly commonly used to compare and combine different PDF sets~\cite{PDF4LHCWorkingGroup:2022cjn}


While any set of basis PDFs can be chosen, the basis of the gluon and individual quark and antiquark flavors (flavor basis) as well as the aforementioned evolution basis are of particular interest. In the NNPDF4.0 determination that we will refer to the basis includes eight independently determined PDFs. Indeed, the charm and anticharm PDFs are assumed to coincide (so  the charm valence sum rule in Eq.~\eqref{eq:val_sumrule} is automatically satisfied), and heavier flavors (bottom and top) are not independently parametrized but rather determined perturbatively. 
The basis PDFs then are, in the flavor basis,
%
\begin{equation}\label{eq:bfldef}
\mathcal{B}_\text{flav}=\{u,\bar u,d,\bar d, s, \bar s, c=\bar c, g\},
\end{equation}
%
and in the evolution basis
%
\begin{equation}\label{eq:bevdef}
\mathcal{B}_\text{ev}=\{V,V_3, V_8,T_3,T_8,T_{15},\Sigma, g\},
\end{equation}
%
where $\Sigma$ is defined in Eq.~(\ref{eq:singlet}), the total valence distribution is
%
\begin{equation}\label{eq:valdef}
V=\sum_{i=1}^{N_f} q^-_i,
\end{equation}
%
and the non-singlet combinations are
%
\begin{align}
T_3&=u^+-d^+, & T_8&=u^++d^+-2 s^+, & T_{15}&=u^++d^+ +s^+-3c^+,\label{eq:tdef}\\
V_3&=u^--d^-, & V_8&=u^-+d^--2 s^-.\label{eq:vdef}
\end{align}


%%------------------------------
%\subsubsection*{FK tables}

%Once the PDFs are parametrized at an initial scale $Q_0$, their remaining degrees of freedom are determined by a global fit to experimental %measurements. 
%Direct numerical evaluation of the double convolution in Eq.\eqref{eq:factorization_theorem} for every data point at every minimization step is %computationally prohibitive. 
%Modern fits therefore rely on precomputed interpolation grids that fold in the DGLAP evolution information, forming so-called \emph{FastKernel %tables (FK tables)}. 
%The convolution integrals are replaced by a discrete tensor contraction on an interpolation grid $\{x_\alpha\}_{\alpha=1}^{N_x}$ at scale $Q_0$,
%%
%\begin{equation}
%\mathcal{O}_n = \sum_{i,j}^{N_{\rm pdf}}\sum_{\alpha,\beta}^{N_x} 
%\text{FK}^{(n)}_{\alpha\beta ij}\;
%f_i(x_\alpha,Q_0)\;
%f_j(x_\beta,Q_0),
%\label{eq:fk_contraction}
%\end{equation}
%
%where $n$ labels the data point, $i,j$ label the PDF combinations in the evolution basis, and $\text{FK}^{(n)}_{\alpha\beta ij}$ encodes the %hard-scattering kernels, bin integrations, and PDF evolution from $Q_0$ to the relevant measurement scale. Once FK tables are available, %evaluating predictions reduces 
%to a set of weighted sums over a fixed $x$-grid, which dramatically accelerates the fit. 

%----------------------------------------
\subsection{The black-box problem}
\label{sec:black_box}

We can view the set of experimental data (about 4000 for NNPDF4.0) as a probability distribution over a high-dimensional data space, typically a multigaussian about the central data points, with width determined by the eigenvalues of the total covariance matrix along its eigenvectors. The regression from data to PDFs is an inverse problem that maps this onto a probability distribution in PDF space. Because of the way predictions are obtained  from PDFs (discussed in Sect.~\ref{sec:factel}), this   inverse problem 
has strong non-localities and correlations. Because  of the way regression is performed in the NNPDF ML-based approach (discussed in Sec.~\ref{sec:pardet}) this leads to the inference chain behaving as a black box, for a variety of reasons that we now summarize.
%
\begin{description}

\item[Discrete data coverage]  Data provide information at a discrete set of points, moreover in a limited kinematic region. The neural net interpolates and extrapolates this discrete information, so that distinct PDF shapes can yield the same value of the loss function. The space of acceptable shapes, and their likelihood, are determined by the hyperparameters of the neural network and its minimization in such a way that the most likely PDF shapes in the final distribution are generally not those leading to the lowest loss, which usually correspond to  overfitted results. 

\item[$x$-region probing] Each observable is sensitive to PDFs that are not evaluated at a single $x$ value, but  rather in the whole region $x_0<x<1$, for some $x_0$ depending on the process kinematics, through the  convolution 
in Eq.\eqref{eq:factorization_theorem}. Hence the interpolation and extrapolation mentioned above actually happen in a conjugate space.

\item[Scale dependence] PDFs are determined at a reference scale $Q_0$, but the  PDFs that enter the physical prediction  Eq.~\eqref{eq:factorization_theorem} are evaluated   at the factorization scale $\mu_F$, in turn chosen of order of a physical scale for each datapoint. The two are related through convolution with an evolution kernel that solves the evolution Eq.~\eqref{eq:evolution}, thereby further modifying in a scale-dependent way  the relation between the $x$ value of the PDF and the kinematics of the data.
   
\item[Scaling violations] Because data are available for several values of scale, and the scale dependence (scaling violation) of the PDF must satisfy the evolution Eq.~\eqref{eq:evolution}, the PDFs are determined not only by the data directly, but also by the scale dependence of the data which is related to the PDF by convolution with the splitting function. Hence the PDF is effectively determined by the mutual relation of data a different scales.


\item[Flavor mixing] Individual quark flavors are a linear combination of different evolution eigenstates, hence they mix with each other as the scale varies, and consequently the linear combinations that appear in a given physical observables are related through this  scale-dependent mixing to the flavor basis PDFs. Even if one chooses the evolution basis, observables cannot be fully  expressed in terms of scale-independent evolution eigenstates because the  singlet and gluon satisfy coupled equations with scale-dependent coefficients.

\item[Non-local correlations] Both data and PDFs that enter different individual predictions of the form of Eq.\eqref{eq:factorization_theorem} are correlated with each other in a highly nonlocal way. The PDFs are correlated because of sum rules 
Eqs.~(\ref{eq:mom_sumrule}-\ref{eq:val_sumrule})  that relate PDFs at all $x$ values. The data are correlated through the covariance matrix $ \Sigma_{nm}^{-1}$ that enters the definition of the loss Eq.~(\ref{eq:loss_simp_pdf}), which may correlate even data coming from independent experiments, but share a common source of theoretical uncertainty.
 
\end{description}
%
Taken together, these effects make it difficult to answer the  simple question,  which PDF flavor, and in which $x$ region, is responsible for the theory prediction to be compared to a specific datapoint. It is important to understand that this is largely a consequence of having approached the inverse problem using machine learning. Indeed, an alternative possibility, adopted by several groups~\cite{PDF4LHCWorkingGroup:2022cjn} is to determine the PDFs by assuming a fixed functional form whose parameters are determined by maximum likelihood. This differs from the machine learning approach firstly, because the dimensionality of the PDF space is rather smaller - typically a few dozens of parameters for a fixed functional form, as opposed to several hundreds for the neural network. And furthermore, because the mapping between the probability distribution in data space and that in PDF space is uniquely fixed by the maximum likelihood condition combined with error propagation, i.e. by the location of the maximum of the likelihood in parameter space, and its behavior in the vicinity of the maximum. In a machine learning approach instead both the most likely PDFs, and its uncertainty are the outcome of a procedure that does not enable a direct determination of the mapping. 


%\clearpage
%%%%%%%%%%%%%%%%%%%%%%%%%%%%%%%%
\section{Shapley values for PDFs}
\label{sec:shapley-values}

Motivated by the black-box structure discussed in Sec.~\ref{sec:black_box}, we approach the problem using a standard XAI idea, namely the Shapley value. In our context, this consists of treating the PDF-to-observable pipeline as a cooperative game, where the players are the individual PDFs and the payoff is the loss. We first briefly  review the game-theoretic foundations of Shapley values and its use as an XAI tool, and then we present our specific implementation in the context of PDF determination,

%---------------------------------------
%\begin{figure}[b!]
%\centering
%\vspace{1pt}
%\input{recap_scheme}
%\caption{Schematic overview of the Shapley value analysis. A baseline PDF set fitted by \texttt{n3fit} is loaded as input; no retraining is performed.
%For a fixed $x$ region and coalition $S$, perturbed PDFs $f^{(S)}$ are propagated through the theory pipeline to obtain predictions and the $\chi^2$. Shapley %values $\phi_j$ quantify the contribution of each flavor for a chosen $x$ region and dataset $\mathcal{D}$.}
%\label{fig:sv4pdf_pipeline}
%\end{figure}
%---------------------------------------

%----------------------------------------
\subsection{The Shapley value in game theory and XAI}

In game theory, the Shapley value is the average marginal contribution $\phi_i$ of a player $i$, determined by comparing the outcomes of a cooperative game 
with and without $i$: $v(S \cup \{i\})$ versus $v(S)$, where $v$ is a value function that maps a coalition, i.e.\ a subset of the total set of $N$ players  $S \subseteq N$ to a real number that gives the payoff achieved by that coalition~\cite{shapley_value1953}. Note that for brevity we denote with $S$ and $N$ both the sets, and the number of elements in the respective sets. 

The Shapley value  is given by
%
\begin{equation}
  \phi_i = \sum_{S \subseteq N \setminus \{i\}} w(S) \left[ v(S \cup \{i\}) - v(S) \right],
  \label{eq:comb-shapley-values}
\end{equation}
%
where the sum runs over all possible coalitions of $N-1$ players not including the $i$-th player, and the weight $w(S)$ ensures that coalitions are weighted accounting for the combinatorial factor that provides the number of different coalitions, namely
the reciprocal of the number of different coalitions with $S$ members. The latter is in turn given by the number of different (non-ordered) subsets of $S$ out of $N-1$ elements, which is the binomial coefficient times the number $N$ of possible choices of the $i$-th element: 
%
\begin{equation}\label{eq:weight}
  w(S) = \left[N\left(\begin{array}{c} N-1\\S\end{array}\right)\right]^{-1}=\frac{S! \, (N - S - 1)!}{N!}.
\end{equation}
%


%---------------------------------------


In machine learning applications, the Shapley value is used in order to determine the contribution of individual features to the loss function: the features are the players, and the loss is the value  function. In this context,  the exact computation of Shapley values is usually intractable, since the typical number $N$ of features is of order of 50 or more, hence the sum over $2^{N-1}$ coalitions is unfeasible. A common way out is the use of the
SHAP framework~\cite{lundberg2017unifiedapproachinterpretingmodel} (Shap Additive Explanation), that addresses this by unifying a class of additive feature attribution 
methods -- including LIME~\cite{ribeiro2016whyitrustyou}, DeepLIFT~\cite{shrikumar2019learningimportantfeaturespropagating}, and LRP~\cite{binder2016layerwiserelevancepropagationneural} -- and 
approximating the Shapley values by assuming feature independence, which then allows for a linear expansion, such that the number of coalitions only grows linearly, rather than exponentially with $N$. In  high-energy physics applications, where input features are typically strongly correlated, this leads to distorted results, as explicitly demonstrated for jet substructure variables in top tagging in Ref.~\cite{Vent:2025ddm}. As discussed in Sect.~\ref{sec:black_box}, in our case PDFs and their features are also strongly correlated, thereby suggesting that this approximation would be equally problematic. As we shall see, we will consider a setup in which the players are the $N_\text{flav}=8$ PDF flavors  $\mathcal{B}_\text{flav}$ or $\mathcal{B}_\text{ev}$ Eqs.~(\ref{eq:bfldef}-\ref{eq:bevdef}). Hence, even though we will consider a large number of different games, in which one flavor is perturbed in a specific region, for each of them we will only have to consider $2^7=128$ coalitions, thereby enabling the computation of the exact Shapley value Eq.~(\ref{eq:comb-shapley-values}).

%----------------------------------------
\subsection{Shapley implementation for PDFs}
\label{sec:setup}

The idea of our approach is to take as a basic feature the behavior of each PDF in a certain $x$ region.  We start with a given set of PDFs determined from a dataset as discussed in Sect.~\ref{sec:pardet}, and we then then consider as a player a perturbed PDF, namely a PDF from this set whose behavior in a given $x$ region has been modified. We then assess the effect of the perturbation by computing the Shapley value for the perturbed PDF in coalition with all other possible coalitions of perturbed PDFs, using as a value function the loss Eq.~(\ref{eq:loss_simp_pdf}) used for PDF determination. 

%\subsubsection{The Shapley value and its uncertainty}

The players are the $N_\text{flav}=8$ flavors in $\mathcal{B}_\text{flav}$ or $\mathcal{B}_\text{ev}$, see Eqs.~\eqref{eq:bfldef}--\eqref{eq:bevdef}. For any subset $S\subseteq\mathcal{B}_\text{flav/ev}$, a coalition
configuration is obtained by perturbing only the flavors in $S$,
%
\begin{equation}\label{eq:coaldef}
    f^{(S)}_j =
        \begin{cases}
          f_j^\text{pert}, & j\in S,\\
          f^{(0)}_j,       & j\notin S,
        \end{cases}
\end{equation}
%
where $f^{(0)}$ denotes the unperturbed baseline PDF set.
The value function is obtained from the loss Eq.~\eqref{eq:loss_simp_pdf}, evaluated using in each instance in the sum over coalitions the PDFs corresponding to coalition $S$ as defined in 
Eq.~(\ref{eq:coaldef}). The perturbation is designed in such a way as  to single out the features that we want to study, namely the behavior of each PDF flavor in a given $x$ region. 

We now discuss the perturbation and corresponding value function in turn. In order to select the behavior in a give $x$ region we 
apply a localized perturbation, chosen to be a  Gaussian bump, thereby ensuring that the smoothness of the unperturbed  PDFs as a function of its argument is preserved.  Because we want to single out the effect of modifying the PDF about its central value in either direction, the perturbation is applied with either a positive or a negative sign. The bump is centered at some $x=x_\mu$,  and has a width $\sigma$ on a logarithmic scale, such  that the PDF is probed on a fixed-width interval on a log scale.  The height of the bump is 
calibrated, namely it is normalized to the PDF uncertainty $\sigma_{\text{PDF},j}(x_\mu)$  at the center of the perturbation, thereby ensuring that the perturbation is physically meaningful and comparable across flavors. We thus have
\begin{align}
\label{eq:perturbation-cal}
    xf_{j}^{\text{pert}\pm}(x,Q^2_0)&=xf_j(x,Q^2_0)+\delta^\pm_j(x),\\\label{eq:deltapm}
        \delta^\pm_j(x)&=\pm \sigma_{\mathrm{rep},j}(x_\mu)
    \exp\!\left[-\frac{\bigl(\log_{10}(x)-\log_{10}(x_\mu)\bigr)^2}{2\sigma^2}\right],
\end{align}
%
A representative example of this calibrated perturbation is shown in Fig.~\ref{fig:10-4calibrated-perturbation}.
%-------------------------
\begin{figure}[t!]
    \centering
    \includegraphics[width=\linewidth]{images/cal_perturbation.pdf}
    \caption{Calibrated perturbation Eq.~(\ref{eq:perturbation-cal}) centered at $x_\mu=10^{-3}$ in
    the flavor basis, with amplitude normalized to $1\sigma$ at $Q_0=1.65~\mathrm{GeV}$ and $\sigma=0.3$.}
    \label{fig:10-4calibrated-perturbation}
\end{figure}
%-------------------------

The value function is then constructed out of the loss in the following way. First, either the perturbation $\delta^+$ or $\delta^-$ Eq.~(\ref{eq:deltapm}) is applied to all PDFs in coalition $S$ Eq.~(\ref{eq:coaldef}). Next, the sum rules Eq.~(\ref{eq:mom_sumrule}-\ref{eq:val_sumrule}) are enforced in the perturbed PDF set: this means that if the coalition $S$ only includes some of the PDFs that enter one of the sum rules, then only these PDFs are perturbed according to Eq.~(\ref{eq:perturbation-cal}), but all the other PDFs entering the sum rules are adjusted by modifying their normalization constants $A_i$ in 
Eq.~(\ref{eq:PDF}) in such a way that the sum rule still holds. This then produces for each coalition two perturbed PDFs $f^{(S+)}$ and $f^{(S-)}$, in which the PDFs in the coalitions are perturbed, and the sum rules always hold. Finally, the value function is evaluated by taking the average of the  loss   Eq.~\eqref{eq:loss_simp_pdf} evaluated for these perturbed PDF sets:
\begin{align}
\label{eq:value-function}
    v(S)=\frac{1}{2}\left[\chi^2\!\left(f^{(S+)};\mathcal{D}\right)+\chi^2\!\left(f^{(S-)};\mathcal{D}\right)\right],
\end{align}
%

The Shapley value is finally computed using this value function in its expression  Eq.~(\ref{eq:comb-shapley-values}), thereby leading to a value $\phi_j(x_\mu)$, depending on the perturbed flavor $j$ and the center $x_\mu$ of the perturbation.  As mentioned, a common way of representing PDF uncertainties, specifically used by NNPDF, is to deliver PDF sets as sets of replicas $f^k$ such that the central expected PDF is the ensemble mean and the uncertainty and correlation are found from the ensemble covariance over the replica sample. Accordingly, we can determine  both the PDF Shapley value and its uncertainty for a given PDF set by evaluating the value $\phi_i$ for each  central PDF $f_i$, and  then  computing the Shapley value  $\phi_i^{(k)}= \phi_j(f^k)$ for each PDF replica $f^k$, i.e.\ evaluating  the value function Eq.~(\ref{eq:value-function}) with fixed data $\mathcal{D}$ and  the $k$-th PDF replica $f_k$. The uncertainty $\sigma_\phi$ is then determined as the  $1\sigma$ standard deviation of $\phi_j$ across the PDF replica ensemble, in turn given by $\sigma_\phi=\sqrt{\text{Var}(\phi_j)}$, with 
\begin{align}
    \bar{\phi}_j &= \frac{1}{N_\text{rep}}\sum_{k=1}^{N_\text{rep}}\phi_j^{(k)},
    \qquad
    \text{Var}(\phi_j) = \frac{1}{N_\text{rep}-1}
    \sum_{k=1}^{N_\text{rep}}\bigl(\phi_j^{(k)}-\bar{\phi}_j\bigr)^2.
    \label{eq:uq}
\end{align}


With our definition, the Shapley value is effectively the marginal change of the value function  Eq.~(\ref{eq:value-function}) when flavor $j$ is added to an existing coalition $S$, with a baseline given by  $v(\emptyset)=\chi^2(f^{(0)};\mathcal{D})$, which is just   the value of the loss which is found when comparing the unperturbed PDF set $f_j$ under investigation to the data $\mathcal{D}$. 
This PDF Shapley value can be computed for a given PDF set using any dataset. This then measures the sensitivity of that PDF set to the given data, which are not necessarily those used in order to determine that PDF set in the first place. Hence this can be used in order to assess the impact that adding new data to some existing PDF determination would have on those PDFs. On the other hand, because it is defined by measuring the impact of a deviation of the given PDF set from its reference best fit, when computed over a subset of  dataset used to determine the given PDF set, it assesses how much of the information on the given PDF set came from that data.

Evaluating the value may lead to three possible situations:
\begin{itemize}
  \item $\phi_j(x_\mu)>0$: the perturbation of flavor $j$ increases $\chi^2$ on average, indicating that the dataset $\mathcal{D}$ is sensitive to and constrains that flavor in the chosen $x$ region.

  \item $\phi_j(x_\mu)<0$: the perturbation improves $\chi^2$ on average, suggesting that the baseline PDF is locally suboptimal for this dataset -- a potential indication of a pull direction in the fit.

  \item $\phi_j(x_\mu)\approx 0$: the dataset is insensitive to flavor $j$ in the chosen $x$ region, either because the flavor does not contribute meaningfully to the relevant predictions or because its effect cancels across coalitions.
\end{itemize}

As a final comment, we note that what allows us to take into account all correlations between features is the small number of coalitions that enables an exact computation of the Shapley index. However, this is only possible because we have defined the coalition of perturbed PDFs according to Eqs.~(\ref{eq:coaldef}), in which the same perturbation is applied to all the PDFs in the coalition so the number of players is equal to $N_{\rm flav}$, and moreover we are conly considering symmetric perturbations according to Eq.~(\ref{eq:value-function}). These are somewhat restrictive assumptions, because, as mentioned, the set of features under investigation is the behavior of each PDF in a certain $x$ region, so in a full analysis, players should be taken as the set of all possible perturbed PDFs, in any $x$ range, and possibly including up and down perturbations. Even if the full physically relevant $x$ range $10^{-4}\lesssim x\lesssim0.7$
was only subdivided in, say four bins for each PDF this would lead to $4 N_{\rm flav}$ players, or twice as many if up and down perturbations are included separately, which corresponds to $2^{32}\sim 10^9$ coalitions, thereby making an exact treatment unmanageable. We shall show in the next section that even within our somewhat simplified setting, the Shapley value performs adequately as a means to probe the impact of data on PDFs.





%
%







%

% In our case, we expect the perturbation to break the sum rules described in Section~\ref{sec:PDF-nutshell}, and for some analysis, it may be interesting for us to allow this unphysical PDF to exists. However, for examples, in cases where the goal of the analysis is to understand the fitting procedure of the neural network, it is interesting to constrain our perturbation to what the NN is constrained to. 
% To ensure Eqs.(~\ref{eq:mom_sumrule},~\ref{eq:val_sumrule}), for each of Shapley coalition we apply a replica of the sum rule layer for each coalition of perturbed PDFs. 
% %
% \begin{align}
%     xf_j^{\text{evol}}(x,Q_0)&=A_j\,x\tilde{f}_{j}(x,Q_0)
%     \\
%     xf_j^{\text{flav}}(x,Q_0)&=A_j\sum_{j'}R_{ji}\,x\tilde{f}_{i}(x,Q_0),
% \end{align}
% %
% where $R_{ji}$ is the rotation matrix. This means that for each coalition we perform correlated perturbation, all PDF of the coalition (perturbed and not perturbed) will be normalized by the sum-rules. Because the number of players is small, we can impose this reconstruction for every coalition and still compute the exact Shapley values. However, this also mean our coalition-wise reconstruction enforces a subset of the physically relevant correlations, namely those induced by the basis transformation and sum-rule constraints, but does not attempt to reproduce the full joint distribution of PDF fluctuations\footnote{A full joint distribution may not be desirable, because we still wish to measure the impact of individual PDFs changes, rather than modifying greatly the full set of PDF for each coalition.}.


%%%%%%%%%%%%%%%%%%%%%%%%%%%%%%%%%
\section{Analyzing the NNPDF4.0 PDF set}
\label{sec:results}

We now provide a first application the technique discussed in the previous section to the NNPDF4.0 PDF set, with the dual goal of checking the viability of our method and exploring its potentialities. We use the default  published  NNLO baseline PDF set
\texttt{NNPDF40\_nnlo\_as\_01180}, with 100 PDF replicas,  and we adopt the same global dataset used  for the determination of these PDFs, as described in Ref.~\cite{NNPDF:2021njg}.  Individual datasets and experiments included in this global datasets are listed in Table~\ref{tab:dataset_full} of  Appendix~\ref{app:datasets}. The kinematic coverage of this dataset in the $(x,Q^2)$ plane is shown in Fig.~2 of Ref.~\cite{NNPDF:2021njg}.

In the NNPDF4.0  determination (unlike in later NNPDF studies) no theory uncertainty contribution associated to missing perturbative higher orders was included. We consequently adopt as a covariance matrix in the value function Eqs.~(\ref{eq:value-function},\ref{eq:loss_simp_pdf}) the full experimental covariance matrix. Note that this differs from the loss used for the training of the neural nets in  Ref.~\cite{NNPDF:2021njg} in  two respects: first, it does not include a contribution introduced in order to account for nuclear uncertainties in datasets involving nuclear targets, and second, it uses the experimental covariance matrix also for the treatment of multiplicative uncertainties, for which a separate uncertainty treatment is necessary in order to avoid biased results. 
Consequently, our value function measures the PDF impact on experimental data calibrated to the purely experimental uncertainty.

We will first evaluate the Shapley value for the full dataset, thus assessing the impact of individual PDFs on the description of the global dataset, thereby providing a first check on the viability of the method. We then turn to individual subsets of data, studying how different PDFs are strongly correlated to different processes in specific kinematic region. While results are by and large in agreement with expectations or previous studies, we find a few mildly surprising results, and one  specific case in which the Shapley index behavior is completely unexpected. Finally, we provide two studies testing the potentialities of our methodology. In the first, we investigate the unexpected results found in the dataset  analysis, and find an likely explanation. Next, we provide a  proof-of-concept application to the design of folds used in the NNPDF hyperparameter optimization procedure.

%----------------------------------------
\subsection{The Shapley value for the  full dataset}
\label{sec:shap_full_data}

We compute the Shapley value for each PDF in turn, using the full global dataset: this provides an estimate of the global impact of PDFs on theory prediction for the NNPDF4.0 dataset.
Before doing this by means of the localized perturbation discussed in Sect.~\ref{sec:setup}, we perform a first check of the method, in which instead of using Eq.~(\ref{eq:perturbation-cal}), we simply set
%
\begin{equation}\label{eq:globzero}
xf_j^{\text{pert}}(x,Q_0^2) = 0.
\end{equation}
%
This cannot be done for the full NNPDF4.0 dataset, as setting some PDFs to zero may lead to the vanishing of some theory predictions, thereby leading to ill-defined values for some observables, typically cross-section ratios. The datasets that have been excluded from the total NNPDF4.0 dataset in this case are listed in Table~\ref{tab:discarded_datasets_ablation}.

%---------------------------
\begin{figure}[h!]
    \centering
    \includegraphics[width=0.99\linewidth]{images/ablation.pdf}\\
    \includegraphics[width=0.95\linewidth]{images/ablation_ev.pdf}
    \caption{Shapley value obtained setting  each PDF to zero in turn according to Eq.~(\ref{eq:globzero}) in the flavor basis (top) and evolution basis (bottom). The dashed horizontal line corresponds to the average over flavors.}
    \label{fig:ablation_sv}
\end{figure}
%---------------------------


The Shapley values are shown in Fig.~\ref{fig:ablation_sv} for each PDF in the flavor (top) or evolution (bottmo) bases, Eqs.~(\ref{eq:bfldef}),~(\ref{eq:bevdef}), respectively.  Comparing the values gives an indication of which PDF has the biggest impact on the theory predictions for the NNPDF4.1 dataset, and conversely, which PDF is mostly constraiuned by the NNPDF4.0 data. Note that the loss function used in the computation of the Shapley value is normalized by the experimental uncertainty in the basis of eigenvectors of the experimental covariance matrix, but the perturbation applied to the PDFs is normalized to the PDF  uncertainty. So  what the Shapley value is telling us is the marginal contribution of that PDF to the theory prediction in units of the data uncertainty. 
We find that the flavor-basis Shapley values are ordered according to
%
\begin{equation}
    \phi_u > \phi_d > \phi_g > \phi_{\bar{u}} > \phi_{\bar{d}} > \phi_s >  \phi_{\bar{s}} > \phi_{c^+},
\end{equation}
%
in agreement with physical expectations: the up and down quarks, whose valence distributions $u^+$ and $d^+$ carry the proton baryon number (see 
Eq.~(\ref{eq:val_sumrule}))  have the largest impact; the gluon, though carrying the largest fraction of the proton momentum, has a somewhat lower impact because its contribution to the majority of processes in the NNPDF4.0 dataset starts at second (next-to-leading) perturbative order; and heavier quarks have an impact that scales as their reciprocal mass.

%---------------------------------------
%\subsubsection{Kinematic dependence}\label{kinematicdependence}
\begin{figure}[b!]
    \centering
    \includegraphics[width=0.95\linewidth]{images/xrange_dense.pdf}\\
    \includegraphics[width=0.95\linewidth]{images/xrange_dense_ev.pdf}
    \caption{Scan of flavor-basis (top) and evolution-basis (bottom) Shapley values as a function of perturbation center $x_\mu$ for NNPDF4.0, computed for the central PDF set. The value $\phi_j(x_\mu)$ is shown for all PDF flavors $j$ and for 50 values of $x_\mu$, equally spaced on a logarithmic scale in the range $ 10^{-4}\le x_\mu\le 0.9 $. Each column corresponds to one value of $x_\mu$, each row to a parton flavor, and the color encodes the signed contribution magnitude on a symlog scale (green: positive, red: negative). The top row shows the mean absolute contribution across flavors, highlighting the overall sensitivity trend with $x_\mu$.}
    \label{fig:dense-x-range}
\end{figure}


After this initial test, we turn to the perturbation Eq.~(\ref{eq:perturbation-cal}). We perform a scan across $50$ logarithmically spaced $x_\mu$ values in the interval 
$[10^{-4}, 0.9]$, with $\sigma= \frac{\log_{10}(0.9)-\log_{10}(10^{-4})}{50}$. Results for the Shapley index computed using the central PDF set are shown as heat maps
in Fig.~\ref{fig:dense-x-range} for all values of the perturbation center, both in the flavor and evolution basis. The values and uncertainties obtained as a mean and standard deviation over the replica sample are shown for six different values of the perturbation center (offset to improve readability) in Fig~\ref{fig:xrange-nnpdf4.0-pos}, where for $x_\mu=0.5$ we also show the distribution of values over replicas for each flavor as a so-called violin plot. 

In order to understand the observed behavior, it is useful to keep in mind that the data are concentrated in the $10^{-3}\lesssim x\lesssim0.5$ region, with the vast majority of data at a scale $Q^2\gtrsim 100$~GeV$^2$ rather higher than the parametrization scale $Q_0$, and that the form of the evolution kernels in Eq.~(\ref{eq:evolution})  is such that the behavior of the initial condition $f_i(x,Q_0^2)$ in some $x\sim x_0$ region mostly affects the behavior of the evolved PDFs $f_i(x,Q^2)$ for $x< x_0$ when the evolved scale is greater than the starting scale $Q>Q_0$. Finally, it should be kept in mind that PDFs are constrained to vanish at the endpoint $x=1$,  and must satisfy the sum rules Eqs.~(\ref{eq:mom_sumrule}-\ref{eq:val_sumrule}). 

Based on these considerations, the overall pattern that emerges from Figs.~\ref{fig:dense-x-range}-\ref{fig:xrange-nnpdf4.0-pos} is clear: the Shapley value differs from zero in the data region, and it vanishes for $x\lesssim 10^{-4}$ where there are no data. Several more detailed features are less obvious, but can be understood based on the known dependence on PDFs of the partonic cross-sections that enter the specific processes in the dataset considered here:
\begin{itemize}
    \item While at very small $x$ the Shapley value fluctuates about zero, one may have expected a rapid drop at the edge of the data region, rather than the slow decrease observed. Indeed, the perturbation is calibrated to the PDF uncertainty, hence as the data become more sparse at small $x$, the PDF uncertainty increases, the  perturbation grows stronger, and one may consequently expect the Shapley value to remain approximately constant. However, it should be kept in mind that at small $x$ the effect of perturbative evolution accorting to Eq.~(\ref{eq:evolution}) leads to a rise of the PDF as $x$ decreases and $Q^2$ increases, which is at leading order independent of the boundary condition, and thus largely washes out the dependence of evolved PDFs on the boundary condition for $x\lesssim 10^{-2}$ in the region where data are available~\cite{Ball:1994du}. This explains the gradual reduction of sensitivity signaled by the decrease of the Shapley value.
    \item The Shapley value remains substantial even for $x_\mu=0.8$, despite the fact that there are almost no data there. This can be understood as a combination of  various effects. First, tha bump at $x\sim 0.8$ at the starting scale  will propagate into PDFs at lower $x$ values at the data scale. Furthermore, as mentioned, the perturbation is calibrated so it will become large at large $x$ as the PDF uncertainty grows, and also, the vanishing constrint at $x=1$ forces the PDF uncertainty to shrink to a point as $x\to1$. Finally, because the momentum sum rule is a constraint on the second moment of the PDF, a bump at large $x$ has an enhanced effect on it, and thus leads to sizable deviations in the  whole $x$ range when the normalization is readjusted in order to maintain the sum rule. These observatios are interesting in view of searches of heavy resonances that probe PDFs at large $x$, and thus one may be interested in understaiding to which extent a signal may or may not be hidden in the PDFs~\cite{Carrazza:2019sec}.
    \item While the flavor nonsinglet valence $V_3$ and $V_8$ combinations are well constrained at large $x\sim 0.3$, the corresponding $q^+$ combinations $T_3$ and $T_8$ are not, and become well constrained only for smaller $x\lesssim 0.2$. This can be understood as a consequence of the fact that flavor nonsinglet combinations are mostly probed in charged-current interactions, such as $W$ production and neutrino deep-inelastic scattering, that probe quark-antiquark differences. The triplet combination $T_3$ is mostly probed by deep-inelastic and fixed-target Drell-Yan data for production on nuclear targets, which are mostly concentrated  at smaller $x$.
    \item The flavor singlet valence combination $V$ is largely unconstrained at large $x$. This mildly surprizing result can be explained by noting that there is in fact no individual data that specifically probes this combination. Specifically, at the LHC, because of the proton-proton initial state, quarks and antiquarks  contribute symmetrically to the contribution to physical observables coming from partonic cross-sections with a quark-antiquark initial state. As it is clear from a comparison of the flavor-basis and evolution basis plots in  Fig.~\ref{fig:dense-x-range}), this  is not apparent when considering individual quark flavors, given that each of them is a combination of singlet and various valence components.  Again this is relevant when looking at signals of new physics, whose quantum numbers may lead to coupling to the valence.
    \item The distributions over replicas shown in Fig.~\ref{fig:xrange-nnpdf4.0-pos} (right) display clear deviations from Gaussian behavior, with an asymmetric shape, a  large number of outliers, concentrated in the upper direction, and possibly even a double-peak shape in the case of the gluon. A similar (though somewhat less pronounced) behavior is observed for the other large $x$ values of the perturbation center, $x_\mu=0.1,\,0.8$, but is completely absent for smaller values of $x_\mu$ (not shown here). This can be understood by noting that in the large $x$ region the momentum sum rule constraint has a large impact because of the $x$ weight in Eq.~(\ref{eq:mom_sumrule}), though less so at $x=0.8$ where the PDF is very small. The momentum sum rule is a nonlocal constraint, hence the effect of the perturbation may affect the prediction for a large number of precisely measured datapoints, along some direction in   data space. The data replicas to which PDF replicas are fitted are generated as multigaussian along directions of eigenvectors of the experimental covariance matrix, but, because of the nonlinear relation between data and observables, the shifts are nonlinearly related to a shift along the direction of the perturbation. Consequently, the effect of the perturbation on individual PDF replicas can be related in  a highly nonlinear way to their shifts, thereby leading to a nongaussian distribution of Shapley values.
\end{itemize}

One particular feature seems surprising and at first puzzling. Namely, the gluon Shapley value, as expected, is large at small $x$ and becomes negligible at very large $x$, in keeping with the expectation that the gluon PDF at the scale of the data is mostly concentrated at low $x$. However, it also  has a very pronounced localized dip between $10^-2$ and $10^-1$, while being sizable both for larger and smaller $x$ values, which is unexpected and somewhat surprising.   We will discuss this behavior and its likely explanation  in Sect.~\ref{sec:gludip} below.








\begin{figure}[t]
    \centering
    \includegraphics[width=1\linewidth]{images/nnpdf40.pdf}
    \caption{Left: The Shapley value and its uncertainty $\phi_j(x_\mu)\pm \sigma_\phi$  obtained from the mean and standard deviation over replicas 
    Eq.~(\ref{eq:uq}), plotted with a symlog scale on the $y$ axis for six different values of the perturbation center $x_\mu$/ The values of  $x_\mu$ are equally spaced logarithmically, and plotted on a logarithmic scale, for all PDF flavors in the flavor basis. The values corresponding to each $x_\mu$ are offset for readability, and are all evaluated at the center of the displayed $x$ bin.  Right: The values $\phi_j(0.5)$ (same as in the  left plot), now shown for all replicas. The width of the colored band corresponds to the density of the replica distribution. The solid three horizontal lines correspond to the mean, maximum, and minimum over the replica sample. Note that the  scale on the $y$ axis is now linear.
    }
    \label{fig:xrange-nnpdf4.0-pos}
\end{figure}


%---------------------------------------
\subsection{Dataset impact}

After testing the behavior of the Shapley value over the full dataset, we now turn to its primary use as a tool in order to study the impact of individual data on PDF determination   and conversely of individual PDFs on specific data predictions. To this purpose, we compute the Shapley value for some specific subset of the NNPDF4.0 full dataset. A complementary exercise would be to compute the Shapley value for data not included in the NNPDF4.0 dataset, but this will be left for future studies. We specifically  study Deep-inelastic scattering (DIS) data, Drell-Yan, i.e. gauge boson  production data, and the combination of jet and top pair production data and in each case, we compute the Shapley for all PDF flavors and for eight values of $x_\mu$ and look at the analogue of Fig.~\ref{fig:xrange-nnpdf4.0-pos}.

%-----------------------------
\subsubsection{Deep-inelastic scattering data}

\begin{figure}[ht]
    \centering
    \includegraphics[width=1\linewidth]{images/DIS.pdf}
    \caption{
    Comparison of flavor-basis Shapley analysis scans across two dataset groups. Panel \textbf{(a)} display the study perform on DIS-HERA datasets, and panel \textbf{(b)} the one performed on DIS fixed target datasets. Both panels show Shapley values for each parton flavor for each calibrated perturbation centered on $x_\mu$. Markers (with error bars) denote flavor-specific central values and uncertainties, with a small horizontal offset applied for readability. 
    }
    \label{fig:DIS}
\end{figure}
Deep-inelastic scattering (DIS) data in principle allow a full determination of all PDFs if both charged-current (neutrino) and neutral-current (charged lepton) data are available, but are only available in a more restricted kinematic region. The Shapley value 

and ,  and indeed have been albeit with bigger  uncertainties and in a more restrictive kinematic region, Drell-Yan data
Deep-inelastic scattering data 
Comparing Fig.~\ref{fig:xrange-nnpdf4.0-pos} and Fig.~\ref{fig:DIS}\textbf{(a)}, we observed for low $x$, $x_\mu \in \{10^{-4},\; 10^{-3}\}$, very similar trend in terms of flavor domination. However, in terms of values, at $x_\mu = 10^{-4}$, the $\phi_j$ are globally greater. In addition, while for the evaluation on NNPDF4.0 like dataset, the values were increasing until $x_\mu=0.5$, the values here never increases.
This trend is what we would expect as DIS-Hera populate low $x$-range observables, meaning that in the NNPDF4.0, they are already greatly responsible for low $x$ region, but at higher $x$, even with completely irrelevant PDF, the impact on the $\chi^2$ will be limited as no points are presents. Impact on the gluon remain limited.

In Fig.~\ref{fig:DIS}\textbf{(b)}, one clear aspect is the low impact of the gluon all over the $x$-range, {\color{red} Should we add more ?}

\begin{figure}
    \centering
    \includegraphics[width=1\linewidth]{images/CC_NC.pdf}
    \caption{
    Comparison of evolution-basis Shapley analysis scans across two dataset groups. Panel \textbf{(a)} displays the study on DIS charged-current (CC) datasets, and panel \textbf{(b)} the one performed on DIS neutral-current (NC) datasets. Both panels show Shapley values for each parton flavor for each calibrated perturbation centered on $x_\mu$. Markers (with error bars) denote flavor-specific central values and uncertainties, with a small horizontal offset applied for readability. 
    }
    \label{fig:CC-vs-NC}
\end{figure}

\Eva{Do we also want to put Fig.~\ref{fig:DIS} there in the evolution basis? Then we could add something like this: 

The HERA collider operates at high COM energies, which results in high possible values of the hard scale of the interaction $Q^2$, which goes up to $10^4$ GeV$^2$. Since this energy is well above the mass of the $Z$ boson, the neutral current interaction takes place and this process is dominated by total quark exchange. Furthermore, the bulk of the HERA data is in the small-$x$ regime, where total quark combinations are present, while valence combinations are absent. This makes HERA data a suitable probe of total quark PDFs. On the contrary, fixed-target datasets are in the larger $x$ regime, and at much lower energies that can be below the mass of the $Z$ boson. In this case, the neutral-current interaction is suppressed, while valence quark combinations become accessible because we are at larger $x$.}

We can also divide the DIS fixed-target datasets into two subsets: the charged-current (CC) datasets obtained from neutrino DIS experiments, and the neutral-current (NC) datasets from electron and muon DIS experiments. We look at the result from this comparison in the evolution basis (see Fig.~\ref{fig:CC-vs-NC}), since in this basis the difference can be directly seen.  For fixed-target NC DIS, the data mainly consists of measurements of the structure function $F_2$. Since $F_2$ depends on total quark combinations $q+\bar{q}$, NC DIS observables can be used to probe the total quark PDFs in the evolution basis. Unlike NC DIS, CC DIS mostly constrains the valence combination of quarks $q-\bar{q}$, due to the asymmetry of the $W$ boson. This makes CC DIS a suitable probe for valence PDFs. The Shapley Value analysis shows this: when performed on fixed-target NC DIS datasets, Shapley Values are nonzero only for $T_3, T_8$ and $\Sigma$. The same analysis on CC DIS datasets shows nonzero Shapley Values also for $V, V_3$ and $V_8$, while the singlet Shapley Value vanishes.


%-----------------------------
\subsubsection{DY datasets}

We separate the total Drell-Yan dataset of distributions into two subsets: Fig.~\ref{fig:DY} shows the Shapley values for a subset with only $p_T$ distributions on the left, and only rapidity distributions on the right. We separate into these two subsets because they are complementary observables. $p_T$ observables depend on the transverse region, while rapidity distributions show the longitudinal behavior. 
The most noticeable difference in Shapley value is the gluon at large-$x$. In the case of the $p_T$ datasets, the gluon PDF clearly has a large contribution to the $\chi^2$, while for the rapidity distributions, $\phi_g$ is very small. This can be explained by the simple argument that for rapidity distributions, the gluon can only appear at higher orders, while for the $p_T$ distributions the gluon is already present in the leading process. This is the case because in order to have a transverse momentum, the interaction needs a recoiling particle, so that the leading order diagrams have either an extra quark or an extra gluon with respect to the leading order diagrams for the rapidity. This makes transverse momentum interactions more sensitive to the gluon PDF.

One effect that is interesting to note is that for the rapidity distributions, at $x_\mu\sim10^{-3}$, the gluon PDF does have a nonzero Shapley Value that is much larger than those of the quarks. In particular, $\phi_g=0.02805$, while for example $\phi_d=0.00893$. This suggests that at this value of $x$, rapidity distributions do have an effect on the gluon PDF, even if it is much smaller than for the $p_T$ distributions.

\begin{figure}[ht]
    \centering
    \includegraphics[width=1\linewidth]{images/DY.pdf}
    \caption{
    Comparison of flavor-basis Shapley analysis scans across two dataset groups. Panel \textbf{(a)} display the study perform on DY $P_T$ datasets, and panel \textbf{(b)} the one performed on DY rapidity datasets. Both panels show Shapley values for each parton flavor for each calibrated perturbation centered on $x_\mu$. Markers (with error bars) denote flavor-specific central values and uncertainties, with a small horizontal offset applied for readability. 
    }
    \label{fig:DY}
\end{figure}


%-----------------------------
\subsubsection{Gluon direct impact datasets}

Jets constrain the gluon because the cross-section depends directly on it. In addition, processes including gluons ($gg \rightarrow gg, qg \rightarrow qg$) are predominant for Jet production at LHC. This directly impact the region from $x\sim 0.02 - 0.5$ due to the energies of the produced jets.
Another process to dominated by gluon is the top pair production $gg\rightarrow t\bar t$, probing in overlapping scale.


\begin{figure}[ht]
    \centering
    \includegraphics[width=1\linewidth]{images/gluon.pdf}
    \caption{
    Comparison of flavor-basis Shapley analysis scans across two dataset groups. Panel \textbf{(a)} display the study perform on Jet datasets, and panel \textbf{(b)} the one performed on $t\bar t$ datasets. Both panels show Shapley values for each parton flavor for each calibrated perturbation centered on $x_\mu$. Markers (with error bars) denote flavor-specific central values and uncertainties, with a small horizontal offset applied for readability. 
    }
    \label{fig:gluon-direct}
\end{figure}
In Fig.~\ref{fig:gluon-direct}, we represent the Shapley values across the $x$-range for each flavor using whether only the jet datasets (see Table~\ref{tab:jet}) and then only the $t\bar t$ datasets (see Table~\ref{tab:tt}). 

For the jet datasets Fig.~\ref{fig:gluon-direct}\textbf{(a)}, we can see an over-domination of the gluon importance at both $x_\mu =0.5$ and  $x_\mu=0.8$. At $x_\mu=0.5$ the $\phi_g=24.10\pm 0.91$, and at $x_\mu=0.8$  $\phi_g=0.35\pm0.02$, while for the rest of the $x$, $\phi_j<0.02$ for any flavor $j$. 

For the $t\bar t$ datasets Fig.~\ref{fig:gluon-direct}\textbf{(b)}, we observe large Shapley values, specifically $\phi_g$ for $x_\mu \in \{0.1,0.5\}$. While remaining smaller than for the jet, it reaches at $x_\mu=0.5$, $\phi_g=0.50\pm 0.04$. Some marginal effect are observed at others $x$. It is interesting to note is at $x_\mu=0.8$, most Shapley values are negative, meaning that for these specific datasets, the $\chi^2$ consistently improved with the perturbation, signaling a weak constraint for this region. This can be explained by the fact that the global NNPDF4.0 fit may put in competition different dataset that better probe this region. 

%----------------------------------------
\subsection{Examples of applications}


\subsubsection{Understanding  gluon determination at intermediate $x$.}
\label{sec:gludip}

In Sec.~\ref{sec:shap_full_data} it was already noted that the Shapley value of the gluon undergoes a sharp dip around $x_\mu\approx 10^{-1}$. We can clearly see this in Fig.~\ref{fig:gluon-dense}, where we observe that $\phi_g$ oscillates around zero between $x_\mu=10^{-1}$ and $x_\mu=2 \times 10^{-1}$, while $|\phi|$, representing the average contribution across all flavors of the indicated basis, remains consistent. In addition, we observe that the Shapley values are consistent between the two bases for all values of $x$.

\begin{figure}[b!]
    \centering
    \includegraphics[width=0.8\linewidth]{images/g_Hera.pdf}
    \caption{Comparison of Shapley values for the gluon PDF in flavor and evolution basis across the kinematic range, evaluated on HERA datasets against the $|\phi|$ average of the Shapley values in all flavors.}
    \label{fig:gluon-dense}
\end{figure}

While the behavior of $\phi_g$ may look unexpected, it actually highlights a feature of QCD evolution. At small-$x$ the presence of the gluon PDF in the nucleon is mainly due to DGLAP evolution, which in Mellin space is given by
%
\begin{equation}
    \frac{\partial g(N,\mu)}{\partial\log \mu} = \frac{\alpha_s(\mu)}{2\pi}\big(\gamma_{qg}(N,\alpha_s(\mu)) + \gamma_{gg}(N,\alpha_s(\mu)\big)g(N,\mu),
\end{equation}
%
where the main contribution comes from the gluon-gluon anomalous dimension $\gamma_{gg}(N,\alpha_s(\mu))$. It is a known fact that the gluon-gluon anomalous dimension crosses zero around $N=2$, causing gluon evolution to be very small at this value of the Mellin variable. Therefore, at this point in Mellin space the gluon PDF is not well constrained by QCD evolution. In order to determine the point in $x$-space at which this zero appears, we can relate the momentum fraction $x$ to the Mellin variable $N$ using a saddle-point argument. The cross section in $x$-space is related to the cross section in $N$-space by the inverse Mellin transform
%
\begin{equation}
    \sigma(x) = \frac{1}{2\pi i}\int_{c-i\infty}^{c+i\infty} \dd N\, x^{-N}\,\sigma(N) = \frac{1}{2\pi i}\int_{c-i\infty}^{c+i\infty} \dd N \,\exp{E(x,N)},
\end{equation}
%
where $E(x,N)\equiv -N\log(x) + \log(\sigma(N))$. The integral can be evaluated using a saddle-point approximation by stating the saddle point condition $\frac{\dd E(x,N)}{\dd N}=0$, which leads to
%
\begin{equation}\label{eq:saddle-point}
\begin{aligned}
    \log(x)&=\frac{\dd\log\sigma(N)}{\dd N}\approx \frac{\dd}{\dd N}\log\big( C_g(N,\alpha_s(Q))g(N,Q)+C_q(N,\alpha_s(Q))q(N,Q) \big),
\end{aligned}
\end{equation}
% 
where in the second equality we write the factorized cross section in terms of the dominating contributions at $\mathcal{O}(\alpha_s)$. $C_i$ denote the coefficient functions corresponding to channel $i$, while $g$ and $q$ are the corresponding PDFs. In the case of gluon evolution at small-$x$ we are mainly interested in the DIS process, since the gluon PDF is mostly constrained by DIS data in this region. We can therefore fill the $\mathcal{O}(\alpha_s)$ $N$-space DIS coefficient functions $C_g(N,\alpha_s(Q))$ and $C_q(N,\alpha_s(Q))$ \Eva{(add in appendix)} into Eq.~\eqref{eq:saddle-point} in order to obtain a relation between $x$ and $N$. For the PDFs we use \texttt{NNPDF40\_nnlo\_as\_01180}. In Fig.~\ref{fig:saddle-point-analysis} we show $x(N)$ for four different energy scales within the range of energy scales observed at HERA. From the figure we see that $N=2$ roughly corresponds to $x=10^{-1}-2 \times 10^{-1}$ for all values of the hard scale $Q$. The dip in the gluon Shapley value around $x_\mu\approx 10^{-1}$ thus corresponds to a dip around $N=2$, which is exactly the point where the gluon anomalous dimension crosses zero. This makes it very likely that the zero Shapley value is indeed due to the fact that gluon evolution tends to zero, which leads to a lack of constraint on the gluon PDF, at that point.

\begin{figure}
    \centering
    \includegraphics[width=0.6\linewidth]{images/saddle-point-analysis.pdf}
    \caption{The momentum fraction $\tau(N)$ as a function of the Mellin variable $N$ from the saddle-point approximation for DIS. The different curves correspond to different values of the hard energy scale of the process, taken to be a range of the HERA energy scales.}
    \label{fig:saddle-point-analysis}
\end{figure}


\subsubsection{Application to fold design}\label{sec:folds}
\begin{figure}[htbp]
    \centering
    \includegraphics[width=0.95\linewidth]{images/gluon_flavor_kfold.pdf}\\
    \includegraphics[width=0.95\linewidth]{images/gluon_evolution_kfold.pdf}
    \caption{}
    \label{fig:gluon_kfold}
\end{figure}

\begin{figure}[p]
    \centering
    \begin{adjustbox}{max totalsize={\textwidth}{0.9\textheight},center}
    \begin{tabular}{cc}
        \includegraphics[width=0.48\linewidth]{images/Sigma_kfold.pdf} &
        \includegraphics[width=0.48\linewidth]{images/V_kfold.pdf} \\

        \includegraphics[width=0.48\linewidth]{images/V3_kfold.pdf} &
        \includegraphics[width=0.48\linewidth]{images/V8_kfold.pdf} \\

        \includegraphics[width=0.48\linewidth]{images/V15_kfold.pdf} &
        \includegraphics[width=0.48\linewidth]{images/T3_kfold.pdf} \\

        \includegraphics[width=0.48\linewidth]{images/T8_kfold.pdf} &
        \includegraphics[width=0.48\linewidth]{images/T15_kfold.pdf}
    \end{tabular}
    \end{adjustbox}
    \caption{Folding: evolution basis}
    \label{fig:evolution_kfold}
\end{figure}

\begin{figure}[p]
    \centering
    \begin{adjustbox}{max totalsize={\textwidth}{0.9\textheight},center}
    \begin{tabular}{cc}
        \includegraphics[width=0.48\linewidth]{images/u_kfold.pdf} &
        \includegraphics[width=0.48\linewidth]{images/ubar_kfold.pdf} \\

        \includegraphics[width=0.48\linewidth]{images/d_kfold.pdf} &
        \includegraphics[width=0.48\linewidth]{images/dbar_kfold.pdf} \\

        \includegraphics[width=0.48\linewidth]{images/s_kfold.pdf} &
        \includegraphics[width=0.48\linewidth]{images/sbar_kfold.pdf} \\

        \includegraphics[width=0.48\linewidth]{images/c_kfold.pdf} &
        \includegraphics[width=0.48\linewidth]{images/cbar_kfold.pdf}
    \end{tabular}
    \end{adjustbox}
    \caption{Folding: flavor basis}
    \label{fig:flavor_kfold}
\end{figure}


\clearpage
%%%%%%%%%%%%%%%%%%%%%%%%%%%%%%%%
\section{Conclusion and outlook}
\label{sec:outlook}

\begin{itemize}
    \item Open possibilities for quantitative combination of dataset folding 
    \item Uncertainty estimation ?
\end{itemize}

In this work, we introduced a novel, model-agnostic approach to interpret the complex pipeline of global PDF determinations using exact Shapley values. By treating PDF flavors as players in a cooperative game and the resulting $\chi^2$ as the value function, we bypassed the opacity inherent in deep-learning-based fits, DGLAP evolution, and cross-dataset correlations. Crucially, the small number of PDF flavors allowed us to compute exact Shapley values, avoiding the feature-independence approximations that frequently limit standard explainable AI methods in highly correlated physical systems.

Applying this methodology to the NNPDF4.0 baseline, we successfully recovered the known hierarchy of PDF flavors and isolated the specific kinematic regions and datasets driving the fit constraints. We demonstrated how distinct observables -- such as DIS, Drell-Yan, and jet production -- exert localized, flavor-specific pulls that perfectly align with perturbative QCD expectations.

Looking forward, this framework opens several promising avenues for PDF phenomenology:
\begin{itemize}
    
\item Quantitative Dataset Folding: Shapley values provide a rigorous metric to assess the true impact of a dataset. By tracking average marginal $\chi^2$ contributions, one could systematically identify redundant measurements or highly impactful data prior to incorporating them into global fits, optimizing the dataset selection process.
\item Uncertainty Estimation and Fit Diagnostics: The distribution of Shapley values across network replicas offers a powerful new lens for uncertainty quantification. Variations in Shapley values can serve as a diagnostic tool to pinpoint under-constrained kinematic regions, over-confident uncertainties, or localized tensions between competing datasets.
\end{itemize}

Ultimately, casting PDF fits into a cooperative game yields a transparent, theoretically grounded tool for interpreting the high-precision constraints required for the next generation of LHC physics.

%%%%%%%%%%%%%%%%%%%%%%%%%%%%%%%%
\section*{Acknowledgements}

Part of this work was supported by from the European Union’s Horizon Europe research and innovation programme under the Marie Sk\l{}odowska-Curie grant agreement No 101168829, Challenging AI with Challenges from Physics: How to solve fundamental problems in Physics by AI and vice versa (AIPHY).

\clearpage
%%%%%%%%%%%%%%%%%%%%
\appendix


%%%%%%%%%%%%%%%%%%%%%%%%%%%%%%%%%
\section{Dataset selection}
\label{app:datasets}

Tables~\ref{tab:dataset_full}--\ref{tab:tt} list the datasets used and discarded in the various subsets of the analysis.
%------------------------------------
\begin{table}[htbp]
\centering
\small
\begin{tabular}{lll}
\toprule
Experiment & Dataset & Reference \\
\midrule
NMC & \texttt{NMC\_NC\_NOTFIXED\_EM-F2} & \cite{}\\
\bottomrule
\end{tabular}
\caption{The NNPDF4.0 dataset. For each dataset we indicate the experiment, the name with which it will be referred to in following  tables and the original reference. More details, such as the kinematic coverage and the number of datapoints, are provided in Tables~1-5 of Ref.~\cite{NNPDF:2021njg}.}
\label{tab:dataset_full}
\end{table}
%------------------------------------
%------------------------------------
\begin{table}[htbp]
\centering
\small
\begin{tabular}{lll}
\toprule
Experiment & Dataset & Critical FK operation \\
\midrule
NMC & \texttt{NMC\_NC\_NOTFIXED\_EM-F2} & ratio \\
E866 & \texttt{DYE866\_Z0\_800GEV\_DW\_RATIO\_PDXSECRATIO} & ratio \\
D0 & \texttt{D0\_Z0\_1P96TEV\_ZRAP} & ratio \\
ATLAS & \texttt{ATLAS\_TTBAR\_8TEV\_LJ\_DIF\_YT-NORM} & ratio \\
ATLAS & \texttt{ATLAS\_TTBAR\_8TEV\_LJ\_DIF\_YTTBAR-NORM} & ratio \\
ATLAS & \texttt{ATLAS\_TTBAR\_8TEV\_2L\_DIF\_YTTBAR-NORM} & ratio \\
CMS & \texttt{CMS\_TTBAR\_8TEV\_LJ\_DIF\_YTTBAR-NORM} & ratio \\
CMS & \texttt{CMS\_TTBAR\_8TEV\_2L\_DIF\_MTTBAR-YT-NORM} & ratio \\
ATLAS & \texttt{ATLAS\_SINGLETOP\_7TEV\_TCHANNEL-XSEC} & ratio \\
ATLAS & \texttt{ATLAS\_SINGLETOP\_13TEV\_TCHANNEL-XSEC} & ratio \\
ATLAS & \texttt{ATLAS\_SINGLETOP\_7TEV\_T-Y-NORM} & ratio \\
ATLAS & \texttt{ATLAS\_SINGLETOP\_7TEV\_TBAR-Y-NORM} & ratio \\
ATLAS & \texttt{ATLAS\_SINGLETOP\_8TEV\_T-RAP-NORM} & ratio \\
ATLAS & \texttt{ATLAS\_SINGLETOP\_8TEV\_TBAR-RAP-NORM} & ratio \\
CMS & \texttt{CMS\_SINGLETOP\_8TEV\_TCHANNEL-XSEC} & ratio \\
CMS & \texttt{CMS\_SINGLETOP\_13TEV\_TCHANNEL-XSEC} & ratio \\
\midrule
D0 & \texttt{D0\_WPWM\_1P96TEV\_ASY} & asymmetry \\
CMS & \texttt{CMS\_WPWM\_7TEV\_ELECTRON\_ASY} & asymmetry \\
CMS & \texttt{CMS\_WPWM\_7TEV\_MUON\_ASY} & asymmetry \\
\midrule
E906 & \texttt{DYE906\_Z0\_120GEV\_DW\_PDXSECRATIO} & combined \\
\bottomrule
\end{tabular}
\caption{Datasets excluded from the zeroing study, grouped by the critical FK operation that prevents their use: ratio, asymmetry, and combined observables.}
\label{tab:discarded_datasets_ablation}
\end{table}
%------------------------------------

%------------------------------------
\begin{table}[htbp]
\centering
\small
\begin{tabular}{lll}
\toprule
Experiment & Dataset & Observable \\
\midrule
HERA & \texttt{HERA\_NC\_318GEV\_EM-SIGMARED} & NC DIS \\
HERA & \texttt{HERA\_NC\_225GEV\_EP-SIGMARED} & NC DIS \\
HERA & \texttt{HERA\_NC\_251GEV\_EP-SIGMARED} & NC DIS \\
HERA & \texttt{HERA\_NC\_300GEV\_EP-SIGMARED} & NC DIS \\
HERA & \texttt{HERA\_NC\_318GEV\_EP-SIGMARED} & NC DIS \\
HERA & \texttt{HERA\_NC\_318GEV\_EAVG\_CHARM-SIGMARED} & charm DIS \\
HERA & \texttt{HERA\_NC\_318GEV\_EAVG\_BOTTOM-SIGMARED} & bottom DIS \\
\bottomrule
\end{tabular}
\caption{Datasets included in the HERA DIS subset.}
\label{tab:HERA}
\end{table}
%------------------------------------


%------------------------------------
\begin{table}[htbp]
\centering
\small
\begin{tabular}{lll}
\toprule
Experiment & Dataset & Observable \\
\midrule
NMC & \texttt{NMC\_NC\_NOTFIXED\_EM-F2} & NC DIS \\
NMC & \texttt{NMC\_NC\_NOTFIXED\_P\_EM-SIGMARED} & NC DIS \\
SLAC & \texttt{SLAC\_NC\_NOTFIXED\_P\_EM-F2} & NC DIS \\
SLAC & \texttt{SLAC\_NC\_NOTFIXED\_D\_EM-F2} & NC DIS \\
BCDMS & \texttt{BCDMS\_NC\_NOTFIXED\_P\_EM-F2} & NC DIS \\
BCDMS & \texttt{BCDMS\_NC\_NOTFIXED\_D\_EM-F2} & NC DIS \\
CHORUS & \texttt{CHORUS\_CC\_NOTFIXED\_PB\_NU-SIGMARED} & CC DIS \\
CHORUS & \texttt{CHORUS\_CC\_NOTFIXED\_PB\_NB-SIGMARED} & CC DIS \\
NUTEV & \texttt{NUTEV\_CC\_NOTFIXED\_FE\_NU-SIGMARED} & CC DIS \\
NUTEV & \texttt{NUTEV\_CC\_NOTFIXED\_FE\_NB-SIGMARED} & CC DIS \\
\bottomrule
\end{tabular}
\caption{Datasets included in the fixed-target DIS subset.}
\label{tab:fixed_target}
\end{table}
%------------------------------------

%------------------------------------
\begin{table}[htbp]
\centering
\small
\begin{tabular}{lll}
\toprule
Experiment & Dataset & Observable \\
\midrule
ATLAS & \texttt{ATLAS\_WJ\_8TEV\_WP-PT} & $W^+$ transverse momentum \\
ATLAS & \texttt{ATLAS\_WJ\_8TEV\_WM-PT} & $W^-$ transverse momentum \\
ATLAS & \texttt{ATLAS\_Z0J\_8TEV\_PT-M} & $Z$ $p_T$ vs invariant mass \\
ATLAS & \texttt{ATLAS\_Z0J\_8TEV\_PT-Y} & $Z$ $p_T$ vs rapidity \\
CMS & \texttt{CMS\_Z0J\_8TEV\_PT-Y} & $Z$ $p_T$ vs rapidity \\
\bottomrule
\end{tabular}
\caption{Datasets included for vector-boson transverse-momentum ($p_T$) distributions in the Drell--Yan sector.}
\label{tab:DY_pt}
\end{table}
%------------------------------------

%------------------------------------
\begin{table}[htbp]
\centering
\small
\begin{tabular}{lll}
\toprule
Experiment & Dataset & Observable \\
\midrule
CDF & \texttt{CDF\_Z0\_1P96TEV\_ZRAP} & $Z$ rapidity \\
D0  & \texttt{D0\_Z0\_1P96TEV\_ZRAP} & $Z$ rapidity \\
ATLAS & \texttt{ATLAS\_WPWM\_7TEV\_36PB\_ETA} & $W^\pm$ pseudorapidity \\
ATLAS & \texttt{ATLAS\_Z0\_7TEV\_36PB\_ETA} & $Z$ pseudorapidity \\
ATLAS & \texttt{ATLAS\_WPWM\_7TEV\_46FB\_CC-ETA} & $W^\pm$ pseudorapidity (central) \\
ATLAS & \texttt{ATLAS\_Z0\_7TEV\_46FB\_CC-Y} & $Z$ rapidity (central-central) \\
ATLAS & \texttt{ATLAS\_Z0\_7TEV\_46FB\_CF-Y} & $Z$ rapidity (central-forward) \\
ATLAS & \texttt{ATLAS\_Z0\_8TEV\_HIMASS\_M-Y} & $Z$ rapidity (high mass) \\
ATLAS & \texttt{ATLAS\_Z0\_8TEV\_LOWMASS\_M-Y} & $Z$ rapidity (low mass) \\
CMS & \texttt{CMS\_WPWM\_8TEV\_MUON\_Y} & $W^\pm$ rapidity \\
LHCb & \texttt{LHCB\_Z0\_7TEV\_DIELECTRON\_Y} & $Z$ rapidity ($e^+e^-$) \\
LHCb & \texttt{LHCB\_Z0\_8TEV\_DIELECTRON\_Y} & $Z$ rapidity ($e^+e^-$) \\
LHCb & \texttt{LHCB\_WPWM\_7TEV\_MUON\_Y} & $W^\pm$ rapidity ($\mu$) \\
LHCb & \texttt{LHCB\_Z0\_7TEV\_MUON\_Y} & $Z$ rapidity ($\mu^+\mu^-$) \\
LHCb & \texttt{LHCB\_WPWM\_8TEV\_MUON\_Y} & $W^\pm$ rapidity ($\mu$) \\
LHCb & \texttt{LHCB\_Z0\_8TEV\_MUON\_Y} & $Z$ rapidity ($\mu^+\mu^-$) \\
LHCb & \texttt{LHCB\_Z0\_13TEV\_DIMUON-Y} & $Z$ rapidity ($\mu^+\mu^-$) \\
LHCb & \texttt{LHCB\_Z0\_13TEV\_DIELECTRON-Y} & $Z$ rapidity ($e^+e^-$) \\
\bottomrule
\end{tabular}
\caption{Datasets included for vector-boson rapidity ($y$) and pseudorapidity ($\eta$) distributions in the Drell--Yan sector.}
\label{tab:DY_rapidity}
\end{table}
%------------------------------------

%------------------------------------
\begin{table}[htbp]
\centering
\small
\begin{tabular}{lll}
\toprule
Experiment & Dataset & Observable \\
\midrule
ATLAS & \texttt{ATLAS\_1JET\_8TEV\_R06\_PTY} & inclusive jets \\
ATLAS & \texttt{ATLAS\_2JET\_7TEV\_R06\_M12Y} & dijets \\
CMS & \texttt{CMS\_2JET\_7TEV\_M12-Y} & dijets \\
CMS & \texttt{CMS\_1JET\_8TEV\_PTY} & inclusive jets \\
\bottomrule
\end{tabular}
\caption{Datasets included in the jet subset.}
\label{tab:jet}
\end{table}
%------------------------------------

%------------------------------------
\begin{table}[htbp]
\centering
\small
\begin{tabular}{lll}
\toprule
Experiment & Dataset & Observable \\
\midrule
ATLAS & \texttt{ATLAS\_TTBAR\_7TEV\_TOT\_X-SEC} & total $t\bar{t}$ \\
ATLAS & \texttt{ATLAS\_TTBAR\_8TEV\_TOT\_X-SEC} & total $t\bar{t}$ \\
ATLAS & \texttt{ATLAS\_TTBAR\_13TEV\_TOT\_X-SEC} & total $t\bar{t}$ \\
CMS & \texttt{CMS\_TTBAR\_5TEV\_TOT\_X-SEC} & total $t\bar{t}$ \\
CMS & \texttt{CMS\_TTBAR\_7TEV\_TOT\_X-SEC} & total $t\bar{t}$ \\
CMS & \texttt{CMS\_TTBAR\_8TEV\_TOT\_X-SEC} & total $t\bar{t}$ \\
CMS & \texttt{CMS\_TTBAR\_13TEV\_TOT\_X-SEC} & total $t\bar{t}$ \\
ATLAS & \texttt{ATLAS\_TTBAR\_8TEV\_LJ\_DIF\_YT-NORM} & $d\sigma/dy_t$ \\
ATLAS & \texttt{ATLAS\_TTBAR\_8TEV\_LJ\_DIF\_YTTBAR-NORM} & $d\sigma/dy_{t\bar{t}}$ \\
ATLAS & \texttt{ATLAS\_TTBAR\_8TEV\_2L\_DIF\_YTTBAR-NORM} & $d\sigma/dy_{t\bar{t}}$ \\
CMS & \texttt{CMS\_TTBAR\_8TEV\_LJ\_DIF\_YTTBAR-NORM} & $d\sigma/dy_{t\bar{t}}$ \\
CMS & \texttt{CMS\_TTBAR\_8TEV\_2L\_DIF\_MTTBAR-YT-NORM} & $d\sigma/dm_{t\bar{t}}$ \\
CMS & \texttt{CMS\_TTBAR\_13TEV\_2L\_DIF\_YT} & $d\sigma/dy_t$ \\
CMS & \texttt{CMS\_TTBAR\_13TEV\_LJ\_2016\_DIF\_YT} & $d\sigma/dy_t$ \\
\bottomrule
\end{tabular}
\caption{Datasets included in the $t\bar{t}$ subset.}
\label{tab:tt}
\end{table}
%------------------------------------

% \clearpage
% %%%%%%%%%%%%%%%%%%%%%%%%%%%%%%%%%
% \section{Perturbation design and uncertainty quantification}
% \label{app:perturbations}

% This appendix details the two perturbation schemes available and the sources of uncertainty entering the Shapley value estimates.

% %--------------------------------------
% \subsubsection*{Perturbation schemes}

% Two perturbation schemes are considered. The \emph{additive} scheme applies a flavor-agnostic Gaussian bump,
% %
% \begin{equation}
% \label{eq:pertubation-add}
%     xf_j^{\text{pert}}(x) = xf_j(x) + \delta(x), \qquad
%     \delta(x) = \varepsilon \exp\!\left(-\frac{(x-\mu)^2}{2\sigma^2}\right),
% \end{equation}
% %
% where $\varepsilon$ is the amplitude, $\mu$ the center, and $\sigma$ the width of the perturbation. While this treats all flavors equally in absolute terms, it is ill-suited to capturing flavor sensitivity when PDF magnitudes differ substantially across flavors: the gluon, for instance, is significantly larger than the quark PDFs in the flavor basis, so an additive perturbation of fixed amplitude has a much weaker relative impact on the gluon.
% The \emph{multiplicative} scheme addresses this by normalizing the perturbation to the local PDF value,
% %
% \begin{equation}
% \label{eq:pertubation-mult}
%     xf_j^{\text{pert}}(x) = xf_j(x)\times\bigl[1 + \delta_j(x)\bigr], \qquad
%     \delta_j(x) = \varepsilon_j \exp\!\left(-\frac{(x-\mu)^2}{2\sigma^2}\right),
% \end{equation}
% %
% ensuring that each flavor receives the same relative perturbation regardless of its absolute magnitude.

% %---------------------
% \begin{figure}[b!]
%     \centering
%     \includegraphics[width=1\linewidth]{images/UQ-compare.pdf}
%     \caption{Shapley value distributions at $\mu=0.5$ under three uncertainty quantification schemes: \textbf{(a)} replica noise only, \textbf{(b)} perturbation sign noise only, and \textbf{(c)} combined replica and sign noise.}
%     \label{fig:UQ-comparison}
% \end{figure}
% %---------------------

% %--------------------------------------
% \subsubsection*{Uncertainty breakdown}

% Two independent sources of uncertainty enter the Shapley value estimates, and it is important to distinguish between them. In all figures, uncertainty brackets represent the $1\sigma$ variance as defined in Eq.\eqref{eq:uq}.
% %
% \begin{itemize}
%     \item The first source is \emph{PDF replica noise}: computing $\phi_j^{(k)}$ separately for each Monte Carlo replica $k$ captures the uncertainty propagated from the experimental data through the fit.

%     \item The second source is \emph{perturbation sign noise}: the sign $\alpha$ entering Eq.\eqref{eq:pertubation-cal} can be drawn randomly from $[-1,+1]$, defining a new game for each sign mask. This would in principle require evaluating $3^9$ coalitions which is computationally prohibitive. The sign-averaged mean and variance are therefore estimated by sampling $N_\text{mask}$ masks via Monte Carlo.
% \end{itemize}
% %
% Both sources are accounted for simultaneously by pairing each of the 100 PDF replicas with a independently drawn sign mask, yielding 100 $(k, \text{mask})$ pairs from which the Shapley value distribution is estimated. Figure~\ref{fig:UQ-comparison} illustrates the effect of each uncertainty source in isolation and in combination.

\clearpage
%%%%%%%%%%%%%%%%%%%%%%%%%%%%%%%%%
\section{Supplementary figures}
\label{app:supplement}

Figures~\ref{fig:ablation_ev}--\ref{fig:dense-x-range} complement the main-text results with evolution-basis Shapley values and a dense $x$-scan in the flavor basis.

%------------------------------------
\begin{figure}[htbp]
    \centering
    \includegraphics[width=0.95\linewidth]{images/ablation_ev.pdf}
    \caption{Shapley values for the ablation perturbation in the evolution basis, complementing the flavor-basis results shown in the main text.}
    \label{fig:ablation_ev}
\end{figure}
%------------------------------------

%------------------------------------
\begin{figure}[htbp]
    \centering
    \includegraphics[width=0.95\linewidth]{images/x-range.pdf}
    \caption{Full NNPDF 4.0 violin flavor basis}
    \label{fig:xrange}
\end{figure}
\begin{figure}[htbp]
    \centering
    \includegraphics[width=0.95\linewidth]{images/x-range-ev.pdf}
    \caption{Evolution-basis Shapley value distributions across NNPDF4.0 datasets. Each panel corresponds to a calibrated perturbation centered on $x_\mu$; each violin shows the replica-level distribution for one parton flavor, with individual replicas overlaid and horizontal markers indicating the mean. The $y$-axis scale varies across panels.}
    \label{fig:xrange-ev}
\end{figure}

\begin{figure}[htbp]
    \centering
    \includegraphics[width=0.95\linewidth]{images/x_range_bar_ev.pdf}
    \caption{Full NNPDF4.0 Evolution basis}
    \label{fig:}
\end{figure}

\begin{figure}[htbp]
    \centering
    \includegraphics[width=0.95\linewidth]{images/x_range_bar_ev.pdf}
    \caption{Full NNPDF4.0 Evolution basis}
    \label{fig:}
\end{figure}

\begin{figure}[htbp]
    \centering
    \includegraphics[width=0.95\linewidth]{images/DIS_ev.pdf}
    \caption{DIS Evolution basis}
    \label{fig:}
\end{figure}
\begin{figure}[htbp]
    \centering
    \includegraphics[width=0.95\linewidth]{images/CC_NC.pdf}
    \caption{cc nc flavor basis}
    \label{fig:}
\end{figure}
\begin{figure}[htbp]
    \centering
    \includegraphics[width=0.95\linewidth]{images/DY_ev.pdf}
    \caption{DY Evolution basis}
    \label{fig:}
\end{figure}

\begin{figure}[htbp]
    \centering
    \includegraphics[width=0.95\linewidth]{images/gluon_ev.pdf}
    \caption{JET TT Evolution basis}
    \label{fig:}
\end{figure}
%------------------------------------

%------------------------------------

%------------------------------------


\clearpage
%%%%%%%%%%%%%%%%%%%%%%

\section{DIS coefficient functions}

In $x$-space, the gluon and quark coefficient functions for DIS at $O(\alpha_s)$ in the $\bar{\text{MS}}$ scheme are given by \cite{Ellis:318585}

\begin{equation}
    \begin{aligned}
        C_q(x) &= \frac{\alpha_s}{2\pi}C_F \Bigg[ 2\left( \frac{\log(1-x)}{1-x} \right)_+ - \frac{3}{2}\left( \frac{1}{1-x} \right)_+ - (1 + x) \log(1-x) \\[4mm]
        &- \frac{1+x^2}{1-x}\log x + 3 + 2x - \left( \frac{\pi^2}{3} + \frac{9}{2} \right) \delta(1-x) \Bigg] \, , \\[4mm]
        C_g(x) &= \frac{\alpha_s}{2\pi}T_R \Bigg[ \left( (1-x)^2 + x^2 \right)\log(\frac{1 - x}{x}) - 8x^2 + 8x - 1 \Bigg] .
    \end{aligned}
\end{equation}

\noindent The corresponding $N$-space coefficient functions are obtained by taking the Mellin transforms of these, and are given by

\begin{equation}
    \begin{aligned}
        \tilde{C}_q(N) &= \frac{\alpha_s}{2\pi}C_F \Bigg[ -\frac{9}{2}+\frac{3}{N} + \frac{2}{N+1} + \frac{3}{2}S_1(N-1) + \frac{S_1(N)}{N} + \frac{S_1(N+1)}{N+1} \\[4mm]
        &- S_2(N-1) - S_2(N+1) + 2 S_{1,1}(N-1)\Bigg], \\[4mm]
        \tilde{C}_g(N) &= \frac{\alpha_s}{2\pi}T_R \Bigg[ \frac{1}{N^2} - \frac{1}{N} + \frac{8}{N+1} - \frac{8}{N+2} - \frac{S_1(N)}{N} \\[4mm]
        &- 2 \left( \frac{1}{(N+1)^2} - \frac{S_1(N+1)}{N+1} \right) + 2 \left( \frac{1}{(N+2)^2} - \frac{S_1(N+2)}{N+2} \right)\Bigg],
    \end{aligned}
\end{equation}

\noindent where $S_i(n)$ are the harmonic sums.



\bibliographystyle{UTPstyle}
\bibliography{ML_ref,PDF_ref}
\end{document}
